\setcounter{secnumdepth}{0}
\section{\centering Анотація}

Кваліфікаційна робота присвячена дослідженню та кількісному оцінюванню якості
україно-англійського автоматичного перекладу усного мовлення. Основну увагу
зосереджено на аналізі семантичної точності, типових помилок перекладу та
міжфразової когерентності.

Як матеріал дослідження використано корпус із 12 сегментів аудіозаписів
політичних звернень українською мовою загальною тривалістю близько 45~хвилин.
Обробка виконана за допомогою каскадного конвеєра, що включає автоматичне
розпізнавання мовлення (ASR, модель Whisper medium) та машинний переклад
(MT, модель OPUS-MT). Додатково проведено порівняльний аналіз із трьома
великими мовними моделями --- GPT-4o, Claude Sonnet 4 та Gemini 2.0 Flash.
Результати оцінено за системою кількісних метрик: WER, CER для ASR; BLEU,
BERTScore-F1, METEOR для MT. Додатково проведено ручну анотацію помилок
(пропуски, спотворення, додавання) та оцінку міжфразової когерентності.

Основні результати: якість ASR --- WER = 8,62\%, CER = 2,10\%; якість MT
(OPUS-MT) --- BLEU = 23,12, BERTScore-F1 = 0,397, METEOR = 0,452; якість MT
(LLM) --- BLEU = 55--57, BERTScore-F1 = 0,71--0,74, METEOR = 0,73--0,75, що
на 62--147\% перевищує показники OPUS-MT ($p < 0,001$). Виявлено статистично
значущу кореляцію між помилками ASR та якістю MT ($r = -0,59$ для CER--BLEU,
$p < 0,05$), що підтверджує каскадний ефект поширення помилок. Ручна анотація
ідентифікувала 114 помилок перекладу OPUS-MT, з яких 79,8\% --- спотворення
(переважно власних назв та географічних назв). LLM-системи демонструють
значно кращу обробку цих категорій завдяки контекстуальним знанням.

\textbf{Ключові слова:} автоматичний переклад мовлення, якість перекладу,
WER, BLEU, BERTScore, каскадний ефект, Whisper, OPUS-MT, великі мовні моделі,
українська мова.
